\documentclass[11pt]{article}
\usepackage{amsmath,amssymb,amsthm,hyperref,geometry}\geometry{margin=1in}
\newtheorem{theorem}{Theorem}\newtheorem{proposition}[theorem]{Proposition}
\title{Residual profiles and admission reflection}
\author{}\date{16 November 2026}
\begin{document}\maketitle
\paragraph{Target and status.} The observable-cover theorem derives assembled readiness once local profiles and exact profile-defined admission are supplied. This note constructs profiles from a declared primitive testing interface and gives an exact existence criterion relative to fixed native witness maps. It separates a genuine impossibility of common-profile presentation from a failure to choose enough profile fields. The construction is a contextual-section quotient, in the tradition of syntactic congruences; no historical novelty or mechanized-verification claim is made.

\section{The primitive testing interface}
Fix a finite component index set and component sets $W_i,W'_i$ and actual native witness maps $r_i:W_i\to W'_i$. These carriers and maps are retained. For each primitive factor $a$, choose a finite scope $S_a$ of components and functions
\[
 F_a:\prod_{i\in S_a}W_i\to V_a,\qquad
 F'_a:\prod_{i\in S_a}W'_i\to V_a.
\]
The functions are the declared tests that consumers may use. Boolean outputs can encode a boundary-matching predicate or raw-admission decision. Other outputs retain an actual event label, task-firing observation, returned abstract interface or observer reading. Partial tests are encoded by a declared undefined value; an unavailable primitive must not silently become a successful test. The declaration identifies the allowed contexts and their information, rather than giving arbitrary actors access to these mathematical tests.

The generated testing contexts here are finite tuples and fixed postprocessings of declared factor outputs, Boolean combinations of Boolean tests, and relational assembly with existential hiding of raw variables. These operations use the declared profiles and their attained domains. Arbitrary operational composition with raw-valued transitions is not included without its own descent or graph-test obligations. Source and target raw admission are defined by the same expression in these factor outputs, for example conjunction of designated Boolean factors. This is local when each factor is genuinely a local boundary/admission operation; renaming a whole-system admission oracle as one primitive would destroy the point. The functions need not have bounded arity, but no efficient finite calculation is promised when scopes are large. A changing or newly admitted consumer adds tests and may require a finer profile.

A common adequate profile presentation consists of $p'_i:W'_i\to P_i$ and $p_i=p'_ir_i$, with equal attained images, and functions $\overline F_a$ on the corresponding attained profile products satisfying
\[
 F'_a=\overline F_a\circ\prod p'_i,\qquad
 F_a=\overline F_a\circ\prod p_i. \tag{P}
\]
Thus the same profile evaluator explains both systems. The condition retains each declared factor, not merely the eventual conjunction of its admission bits. It is intentionally stronger than preserving only one globally reachable outcome.

\section{Constructing the coarsest adequate target profiles}
For $x,y\in W'_i$, define $x\equiv_i y$ when every incident primitive test has equal sections:
\[
 F'_a(z[i:=x])=F'_a(z[i:=y])
 \quad\text{for every }a\text{ with }i\in S_a
 \text{ and all }z\in\prod_{k\in S_a\setminus\{i\}}W'_k.
 \tag{R}
\]
Let $q_i:W'_i\to Q_i=W'_i/{\equiv_i}$ be the quotient. An isolated component has the indiscrete equivalence, unless unary observations were specified. Empty witness sets and empty products are interpreted literally.

\begin{proposition}[Residual quotient and its universal property]\label{prop:residual}
Every $F'_a$ factors uniquely through the products of the $q_i$. If profiles $p'_i$ make every $F'_a$ factor, then $q_i=k_ip'_i$ for a unique $k_i$ on the attained image of $p'_i$. Thus $(q_i)$ is the coarsest adequate family of componentwise profiles, unique up to the canonical quotient isomorphisms.
\end{proposition}
\begin{proof}
Each $\equiv_i$ is an intersection of kernels of section functions and hence an equivalence. To compare two tuples with coordinatewise equivalent entries, replace their finitely many coordinates one at a time; equation (R) preserves $F'_a$ at each replacement. This defines the quotient evaluator, uniquely since the product quotient is surjective for a finite scope. If $p'_i(x)=p'_i(y)$, hold all other raw target arguments fixed. The assumed factorization makes the two test results equal for every incident factor, so $x\equiv_i y$. Therefore $q_i$ is constant on fibres of $p'_i$ and descends uniquely to its image.
\end{proof}

The word coarsest is relative to the displayed primitive tests with all surrounding target arguments admitted as mathematical probes. Restricting probes to reachable or authorized contexts can permit more identification; that is a different test interface. The quotient need not be small. It derives a representation from actual consumers instead of selecting a convenient summary and hoping no consumer distinguishes its fibres.

\section{Exact compatibility with fixed native comparisons}
\begin{theorem}[Existence and construction of common profiles]\label{thm:criterion}
A common adequate profile presentation (P) exists if and only if both conditions hold:
\begin{enumerate}
\item Every primitive test is reflected along the given witness maps:
\[
 F_a=F'_a\circ\prod_{i\in S_a}r_i. \tag{T}
\]
\item Every residual target class has a source representative:
\[
 q_i r_i:W_i\twoheadrightarrow Q_i\quad\text{for every }i. \tag{C}
\]
When these conditions hold, take $P_i=Q_i$, $p'_i=q_i$ and $p_i=q_ir_i$. All factor-defined admission and observations then descend, furnishing the local covers and actual profile-defined assembly required by the observable-cover theorem.
\end{theorem}
\begin{proof}
For sufficiency, Proposition~\ref{prop:residual} constructs $\overline F_a$ for the target. Substitute (T) to obtain the same evaluator on source profiles. Condition (C) gives equal attained images, both $Q_i$.

For necessity, substitute $p_i=p'_ir_i$ in (P), proving (T). Proposition~\ref{prop:residual} gives $q_i=k_ip'_i$. Since $p'_i(W'_i)=p_i(W_i)$, each target profile and therefore each residual class is attained by a source witness, giving (C).

For the assembly conclusion, let $D\subseteq\prod Q_i$ be the relation defined by the chosen expression in quotient factor outputs. A tuple of source witnesses is admitted exactly when its profile tuple belongs to $D$, and likewise for the target. Hence the actual carriers are precisely the pullbacks of $D$ along the component profile products, or equivalently the fibre products in the observable-cover theorem. No global carrier-image equality has been assumed.
\end{proof}

This criterion is local in the factorization: reflection and residual sections concern each declared primitive's scope, not the product of all project components. Residual classes for a component combine the signatures of all incident factors. Cycles in the incidence graph do not affect the proof. In finite tables the criterion is executable. It may still be expensive, and for infinite or dynamically generated carriers the equations need symbolic proofs or a suitable reachable invariant.

Failure of (C) cannot be repaired by choosing \emph{finer} adequate common profiles: every such profile maps onto the residual quotient, so any cover would cover the missing residual class. Failure of (T) is more severe: no choice of common profiles compatible with the same $r_i$ can make the factor commute. This is an impossibility result for that interface, not an impossibility of a different comparison or richer boundary.

\section{Why tree lifting and common profiles solve different problems}
The existing local-joint manuscript proves full joint witness surjectivity on a tree from root surjectivity and relative child surjectivity
\[
 (r_i,b_i):W_i\twoheadrightarrow W'_i\times_{B'_e}B_e.
\]
This lifts target witnesses by choosing compatible source boundary values. It does not say that every tuple of source preimages of a permitted target tuple is itself permitted.

Consider two source bit carriers with equality admission, target singleton carriers with unrestricted admission, and constant maps $r_i$. Collapse the source bit boundary to the singleton target boundary. The source diagonal covers the target singleton. The root map is surjective and the relative child map to $1\times_1\{0,1\}$ is bijective. Nevertheless source equality is not the inverse image of target True: the source tuple $(0,1)$ is rejected. Thus (T) fails for the boundary test. No enrichment of \emph{common profiles compatible with these fixed maps} can represent the source equality, since every such source profile is constant.

Conversely, a source singleton may map into a target with two private witnesses having identical residual signatures. Condition (C) holds while full target witness surjectivity fails. Thus the two certificates are genuinely different, not alternative proofs of one unstated universal theorem.

The distinction is exactly between \emph{existence of some compatible lift} and \emph{reflection of admission for every given source tuple}. Readiness may need only the first in some settings; the pullback presentation of common profiles needs the second. Rejecting a common-profile construction is not evidence that a valid native comparison or the earlier tree certificate is unavailable.

\section{Enriching the interface without losing native witnesses}
When reflection fails because a target boundary erased a source compatibility distinction, a legitimate repair may retain an additional boundary witness on the target side. The earlier relative object $W'_i\times_{B'_e}B_e$ indicates the data needed for a chosen boundary. Replacing $W'_i$ by such an object changes the comparison interface; it is not merely renaming a profile of the old target.

An enriched native comparison still owes a joint effect, actual grade arrows, coherent lifting, and chooser descent across inputs. A source effect can be pushed to an enriched carrier at one fixed input, but this does not automatically descend to a chooser shared by all source inputs with the same target input. The original coefficient-limit and admission/chooser theorems remain relevant. This note neither chooses new native lifts nor turns existential lifts into coherent ones.

The residual quotient is a \emph{certificate interface}, not a replacement for the native witness or coefficient object. It must not identify different native arrows merely because their test outputs agree. Probability and coefficient transport continue to use the actual $r_i$ and their native data. The construction can reveal which distinctions must be exposed for a stated interaction without quotienting away the comparison evidence itself.

\section{Sequential use and changing obligations}
At an actual state and command, instantiate the factors with that input. The theorem supplies one-step readiness data. To use it at every execution prefix, prove (T)/(C) for every admitted input or a preserved invariant; finite testing at selected states is not enough. Returned abstract state, actual event labels and observer readings must be included among tests whenever the next rule or policy consumes them. Dependent continuations must descend through those returned interfaces. The local-observer update theorem then derives history-information factorization separately.

Adding a test refines each residual equivalence. The old target profile remains adequate exactly when the new target test is constant on its profile fibres. The old common certificate remains adequate when, additionally, the corresponding source test satisfies reflection (T), so both tests descend through the same evaluator. This gives an explicit invalidation test: a new position-sensitive fault or observer may destroy a count-only profile. There is no assertion that the first quotient anticipates every future consumer. Unknown consumers remain a reason to retain reserve information rather than silently discard it.

\paragraph{Verification.} The written criterion and its boundary cases were independently reviewed. The finite diagnostic compares the criterion with exhaustive enumeration of adequate profile partitions for all Boolean binary tests with target cardinalities $(3,2)$, source cardinalities $(2,2)$ and all maps between them: 36,864 comparisons, plus coarseness and boundary assertions. This is not exhaustive testing of all finite interfaces and not proof-assistant certification. TeX compilation is not claimed.

\section{Literature and claim boundary}
Equality under all declared one-hole tests is the standard syntactic-congruence pattern. For a primary author-hosted treatment of the classical semigroup version, see Tero Harju, \emph{Lecture Notes on Semigroups}, Section 2.1, Lemma 2.3, \url{https://users.utu.fi/harju/semigroups/semigroups96.pdf}. The many-sorted factor theorem above is proved directly; no language-recognition theorem is imported as a native lifting result.

The image/pullback consequence uses the regular-image calculus already credited in \path{Observable_Covers_and_Local_Observers.tex}. The existing tree criterion is in the raw-coverage section of \path{../../foundations/local-joint-comparison/Local_Joint_Comparison_and_Generated_Interference.tex}. No claim of optimality over all possible theories, probes, state restrictions or comparison spans is made. The exact result concerns the fixed native witness maps and explicitly declared factor tests.
\end{document}
